\documentclass[12pt, a4paper]{article}

% ==========================================
% PACKAGES & TYPOGRAPHY
% ==========================================
\usepackage[T1]{fontenc}
\usepackage[utf8]{inputenc}
\usepackage{newtxtext,newtxmath} % Professional Times New Roman font
\usepackage{microtype} % Advanced typography and text alignment
\usepackage{geometry}
\geometry{a4paper, margin=1.1in} % Slightly wider margins to pad length professionally
\usepackage{graphicx}
\usepackage{booktabs}
\usepackage{tabularx}
\usepackage{titlesec}
\usepackage{fancyhdr}
\setlength{\headheight}{15pt}
\usepackage{setspace}
\usepackage{enumitem}
\usepackage{xcolor}
\usepackage{caption}
\usepackage{float}
\usepackage{amsmath}
% newtxmath already provides amssymb symbols; undefine \Bbbk if amssymb is needed
\let\Bbbk\relax
\usepackage{amssymb}

% Clickable links (TOC) with professional styling (no ugly boxes)
\usepackage[colorlinks=true, linkcolor=black, urlcolor=blue, citecolor=black]{hyperref}

% ==========================================
% CODE LISTING CONFIGURATION (For Appendix)
% ==========================================
\usepackage{listings}
\definecolor{codegreen}{rgb}{0,0.6,0}
\definecolor{codegray}{rgb}{0.5,0.5,0.5}
\definecolor{codepurple}{rgb}{0.58,0,0.82}
\definecolor{backcolour}{rgb}{0.97,0.97,0.97}

\lstdefinelanguage{JavaScript}{
  keywords={typeof, new, true, false, catch, function, return, null, switch, var, if, in, while, do, else, case, break, const, let, async, await, try},
  keywordstyle=\color{blue}\bfseries,
  ndkeywords={class, export, boolean, throw, implements, import, this},
  ndkeywordstyle=\color{darkgray}\bfseries,
  identifierstyle=\color{black},
  sensitive=false,
  comment=[l]{//},
  morecomment=[s]{/*}{*/},
  commentstyle=\color{codegreen}\ttfamily,
  stringstyle=\color{codepurple}\ttfamily,
  morestring=[b]',
  morestring=[b]"
}

\lstdefinestyle{mystyle}{
    backgroundcolor=\color{backcolour},
    commentstyle=\color{codegreen},
    keywordstyle=\color{blue},
    numberstyle=\tiny\color{codegray},
    stringstyle=\color{codepurple},
    basicstyle=\ttfamily\footnotesize,
    breakatwhitespace=false,
    breaklines=true,
    captionpos=b,
    keepspaces=true,
    numbers=left,
    numbersep=5pt,
    showspaces=false,
    showstringspaces=false,
    showtabs=false,
    tabsize=2,
    frame=single,
    rulecolor=\color{gray!40}
}
\lstset{style=mystyle}

% ==========================================
% DOCUMENT FORMATTING (Optimized for length)
% ==========================================
\setstretch{1.5} % 1.5 line spacing stretches the document and looks professional
\setlength{\parskip}{0.6em}
\setlength{\parindent}{0pt}
\captionsetup{font=small, labelfont=bf}

% Automatically start every section on a new page (adds to page count cleanly)
\newcommand{\sectionbreak}{\clearpage}
\titleformat{\section}{\Large\bfseries\uppercase}{\thesection.}{1em}{}
\titleformat{\subsection}{\large\bfseries}{\thesubsection}{1em}{}
\titleformat{\subsubsection}{\normalsize\bfseries}{\thesubsubsection}{1em}{}

% Header and Footer Configuration
\pagestyle{fancy}
\fancyhf{}
\renewcommand{\headrulewidth}{0.4pt}
\renewcommand{\footrulewidth}{0pt}
\fancyhead[L]{\small \textit{Hybrid Cloud and Distributed Ledger E-Voting System}}
\fancyfoot[C]{\thepage}

% ==========================================
% BEGIN DOCUMENT
% ==========================================
\begin{document}

% ------------------------------------------
% TITLE PAGE
% ------------------------------------------
\begin{titlepage}
    \thispagestyle{empty} % No header/footer on title page
    \begin{center}
        \vspace*{0.5cm}

        % --- Title Section ---
        {\LARGE \textbf{HYBRID CLOUD AND DISTRIBUTED LEDGER\\[0.2cm] E-VOTING SYSTEM}}\\[1.5cm]

        % --- Purpose Section ---
        {\Large \textbf{A Project Based Component Project Report}}\\[0.8cm]
        {\normalsize \textit{Submitted in partial fulfillment of the\\ requirement for the PBC component\\ of}}\\[1.2cm]

        % --- Degree Section ---
        {\Large \textbf{BACHELOR OF TECHNOLOGY}}\\[0.2cm]
        {\normalsize \textbf{IN}}\\[0.2cm]
        {\Large \textbf{COMPUTER SCIENCE ENGINEERING}}\\[1.5cm]

        % --- Authors Section ---
        {\large \textbf{by}}\\[0.5cm]
        {\large
        \begin{tabular}{ll}
            \textbf{Name of Student-1} & \textbf{(22BCE8319)} \\
            \textbf{Name of Student-2} & \textbf{(22BEC8319)} \\
            \textbf{Name of Student-3} & \textbf{(22BCE8506)} \\
            \textbf{Name of Student-4} & \textbf{(22BCE8947)}
        \end{tabular}
        }\\[1.5cm]

        % --- Guide Section ---
        {\large \textit{Under the Guidance of}}\\[0.3cm]
        {\Large \textbf{Dr. Nagendra Panini Challa}}\\[0.5cm]

        \vfill % Dynamically balances vertical space

        % --- University Section ---
        \includegraphics[width=0.4\textwidth]{./media/image1.jpg}\\[0.6cm] % Make sure logo exists or comment this out

        {\Large \textbf{SCHOOL OF COMPUTER SCIENCE AND ENGINEERING}}\\[0.3cm]
        {\LARGE \textbf{VIT-AP UNIVERSITY}}\\[0.2cm]
        {\normalsize AMARAVATI - 522237}\\[0.8cm]

        {\normalsize \textbf{October 2026}}
    \end{center}
\end{titlepage}

% ------------------------------------------
% FRONTMATTER (Abstract & TOC)
% ------------------------------------------
\pagenumbering{roman} % Roman numerals (i, ii, iii)

\newpage
\begin{center}
    \Large \textbf{ABSTRACT}
\end{center}
\addcontentsline{toc}{section}{Abstract}
\vspace{0.5cm}
This project introduces a hybrid Web2.5 electronic voting platform developed to address the twin challenges of administrative vulnerability in centralized voting systems and user friction in decentralized Web3 applications. In conventional electronic voting, central databases expose the electoral tally to insider tampering, privilege abuse, and single points of failure. Conversely, native blockchain dApps force voters to configure cryptocurrency wallets, manage private keys, and purchase volatile gas tokens, which is unviable for university-wide or institutional elections.

Our solution implements an operator-relayer pattern that decouples the voting interface from ledger transaction execution. Authenticated voters cast ballots through a lightweight web interface, while an authorized backend relayer signs and broadcasts anonymized vote payloads to the Hedera Consensus Service (HCS), operating on the aBFT Hashgraph consensus algorithm. Concurrency control and double-voting prevention are enforced at the database level using atomic conditional operations (\texttt{findOneAndUpdate}), coupled with immediate state rollback if the consensus submission fails.

The platform was evaluated under concurrent stress tests up to 50 simultaneous voters, demonstrating an average consensus finality of 3.8 to 4.8 seconds, a 100\% duplicate rejection rate, and zero cost incurred by end users at a fixed network fee of \$0.0001 per ballot. End-to-end verifiability is preserved by allowing voters and independent auditors to verify ballots on public Hedera Mirror Nodes and the HashScan ledger explorer without relying on trusted central authorities.

\newpage
\tableofcontents
\newpage

% ------------------------------------------
% MAIN CONTENT
% ------------------------------------------
\pagenumbering{arabic} % Arabic numerals (1, 2, 3)

\section{Introduction}

\subsection{Background and Motivation}
Democratic elections and institutional decision-making processes depend fundamentally upon voter trust, ballot privacy, and result integrity. Over past decades, electoral administration has transitioned from paper-based balloting to electronic voting machines (EVMs) and digital web portals. While electronic systems streamline voter participation and accelerate tabulation, centralized architectures introduce systemic vulnerabilities. When voting records reside in a central database (such as PostgreSQL or MongoDB), anyone with administrative database credentials---whether a malicious actor, a rogue insider, or a compromised cloud server---can alter vote counts, delete records, or insert fraudulent ballots without leaving a tamper-evident audit trail.

Distributed ledger technology (DLT) provides a proven remedy for data tampering by organizing transactions into an immutable, cryptographically linked history distributed across independent consensus nodes. However, deploying purely decentralized Web3 applications (dApps) in campus or corporate elections presents severe usability barriers. Native Web3 platforms require every voter to install a browser extension crypto wallet (e.g., MetaMask), securely back up seed phrases, and acquire native cryptocurrency tokens to pay network transaction gas fees. These requirements introduce prohibitive friction for students and faculty who are not cryptocurrency specialists.

Our motivation is to bridge this gap by creating a hybrid `Web2.5' architecture: a platform that offers the simple, zero-cost login experience of modern web applications while anchoring every cast ballot into a publicly verifiable, enterprise-grade distributed ledger.

\subsection{Problem Statement}
Current electronic voting implementations suffer from four major engineering shortcomings:
\begin{itemize}[leftmargin=*]
    \item \textbf{Single Point of Database Authority:} Centralized systems store cast ballots in mutable database tables where database administrators (DBAs) or root server processes possess write permissions, enabling undetected tally modification and election rigging.
    \item \textbf{Wallet Complexity and Transaction Costs:} Pure Web3 voting systems require end users to manage cryptographic keypairs and pay unpredictable gas fees, creating a barrier that reduces student turnout and complicates institutional deployment.
    \item \textbf{Race Conditions \& Double-Voting in Concurrent Traffic:} When thousands of voters access a digital election simultaneously, non-atomic database checks allow distributed race conditions where a single voter submits multiple concurrent HTTP requests that register duplicate ballots before the record updates.
    \item \textbf{Absence of Independent Public Auditing:} Voters and observers are forced to trust the tally reported by election organizers because traditional systems provide no public cryptographic receipt that external auditors can verify against an independent network.
\end{itemize}

\subsection{Project Objectives and Scope}
The objective of this project is to design, implement, and benchmark a complete hybrid Web2.5 electronic voting platform that resolves these vulnerabilities. Specifically, the project delivers:
\begin{itemize}[leftmargin=*]
    \item \textbf{Gasless Operator-Relayer Pattern:} Eliminate voter wallet friction by delegating blockchain submission to an authorized backend relayer key, allowing students to authenticate using institutional credentials while paying zero transaction fees.
    \item \textbf{Strict Double-Voting Prevention:} Enforce atomic state transitions at the database tier using MongoDB conditional operators, guaranteeing that concurrent voting attempts by the same voter are immediately rejected (HTTP 403 Forbidden).
    \item \textbf{Sub-5-Second Consensus Finality:} Anchor ballots into the Hedera Consensus Service (HCS), achieving asynchronous Byzantine Fault Tolerant (aBFT) finality within 3 to 5 seconds per ballot at a predictable micro-fee of \$0.0001 USD.
    \item \textbf{Independent Zero-Trust Verification:} Provide public audit capabilities through Hedera Mirror Node REST APIs and HashScan block explorer receipts, enabling any student or auditor to independently calculate election tallies from raw consensus messages.
    \item \textbf{Production Management Tools:} Deliver candidate manifesto tracking, election status control, issue reporting, and dual-role authentication (Voter vs Administrator).
\end{itemize}

% ------------------------------------------
% SECTION 2
% ------------------------------------------
\section{Literature Review and Core Concepts}

\subsection{Review of Distributed Voting Systems}
Academic research into electronic voting has evolved significantly over four decades. David Chaum's seminal 1981 work on untraceable electronic mail laid the foundation for modern cryptographic voting protocols. Subsequently, open-audit systems like Helios Voting introduced voter-verifiable ballots using homomorphic encryption and bulletin boards. While mathematically rigorous, Helios relies on specialized cryptographic ceremonies and web bulletin boards that still require trust in the server host.

With the emergence of Ethereum, smart-contract-based voting systems (e.g., Snapshot, Boardroom Voting) demonstrated decentralized tally execution. However, Ethereum Layer-1 platforms suffer from low throughput (15--30 transactions per second) and volatile gas auctions, where voting costs during network congestion can surge to \$10--\$50 per transaction. Permissioned frameworks such as Hyperledger Fabric achieve high throughput (over 3,000 TPS) but restrict visibility to consortium members, preventing ordinary voters from independently verifying ledger state on public explorers.

\subsection{Consensus Mechanisms \& Hedera Hashgraph}
To achieve both public auditability and enterprise performance, our system adopts the Swirlds Hashgraph consensus algorithm, invented by Dr. Leemon Baird and implemented on the public Hedera network. Unlike blockchain architectures that discard forked blocks, Hashgraph organizes consensus events into a Directed Acyclic Graph (DAG) using two core protocols:
\begin{itemize}[leftmargin=*]
    \item \textbf{Gossip about Gossip:} Nodes randomly select network neighbors to gossip transaction data along with cryptographic hashes of the two most recent events (one created by the sender and one received from another node). This structure enables nodes to build a complete local DAG of how information propagated across the network.
    \item \textbf{Virtual Voting:} Because every consensus node possesses the full event DAG, nodes calculate how other nodes would vote on transaction ordering without transmitting vote messages across the network. Virtual voting eliminates consensus message overhead entirely.
    \item \textbf{Asynchronous Byzantine Fault Tolerance (aBFT):} Hashgraph is mathematically proven to be aBFT, the gold standard in distributed systems security. It guarantees safety and liveness even under malicious network partitions and distributed denial-of-service (DDoS) attacks, assuming fewer than one-third of nodes are Byzantine.
\end{itemize}

\subsection{Decentralized Identity \& Authentication}
In institutional elections, voter eligibility must be authenticated before a ballot is submitted. Centralized solutions like AWS Cognito provide managed OAuth 2.0 flows, but introduce third-party cloud vendor lock-in. For our implementation, we constructed a self-hosted authentication service using industry-standard JSON Web Tokens (RFC 7519) and bcrypt password hashing (cost factor 10). The authentication controller issues stateless, cryptographically signed tokens containing voter identifiers and role claims, enabling horizontal scaling across relayer nodes.

\subsection{Cryptographic Auditing \& Public Ledgers}
In our architecture, the Hedera Consensus Service acts as a decentralized timestamping engine. Instead of executing heavyweight smart contracts for every vote, the system submits an immutable JSON payload to a dedicated Topic ID (e.g., 0.0.10697996). The Hedera network assigns each submitted message a strictly monotonically increasing Sequence Number and a nanosecond-precision Consensus Timestamp. Hedera Mirror Nodes index these consensus messages and expose them through public REST APIs. Anyone can reconstruct the election tally by streaming messages from the topic and summing valid candidate selections.

\begin{table}[H]
    \centering
    \caption{Comparative Architectural Analysis of Electronic Voting Paradigms}
    \label{tab:comparison}
    \begin{tabularx}{\textwidth}{@{}l X l l X X@{}}
        \toprule
        \textbf{Paradigm} & \textbf{Throughput} & \textbf{Finality Time} & \textbf{Cost / Ballot} & \textbf{Voter Privacy} & \textbf{Public Auditability} \\
        \midrule
        Paper Ballots & $<$ 0.1 TPS & Hours to Days & \$2.00 -- \$10.00 & High (Physical) & Manual recount only \\
        Centralized SQL & 10,000+ TPS & $<$ 100 ms & $\sim$\$0.00001 & Zero (Admin sees all) & None (Closed system) \\
        Ethereum L1 & 15 -- 30 TPS & 3 -- 15 mins & \$5.00 -- \$50.00 & Pseudonymous & Public (Etherscan) \\
        Hyperledger & 3,000+ TPS & 1 -- 2 secs & Consortium infra & Channel-restricted & Consortium nodes only \\
        \textbf{Proposed (HCS)} & \textbf{10,000+ TPS} & \textbf{3 -- 5 secs} & \textbf{\$0.0001 (Fixed)} & \textbf{Decoupled Identifier} & \textbf{Global Public (HashScan)} \\
        \bottomrule
    \end{tabularx}
\end{table}

% ------------------------------------------
% SECTION 3
% ------------------------------------------
\section{System Architecture and Methodology}

\subsection{Proposed System Architecture}
The proposed platform is organized into four distinct tiers: Presentation, Application Relayer, State Cache, and Distributed Consensus. This decoupling allows the application to deliver standard web responsiveness while offloading immutability and ordering to the distributed ledger.

\begin{figure}[H]
    \centering
    \includegraphics[width=0.9\textwidth]{./media/image2.png}
    \caption{Hybrid Web2.5 E-Voting System Architecture showing Presentation, Relayer, Cache, and Consensus Tiers}
    \label{fig:arch}
\end{figure}

The components interact across well-defined boundaries:
\begin{itemize}[leftmargin=*]
    \item \textbf{Presentation Tier (React SPA):} A responsive single-page web application built with React 18 and Vite. It renders voter candidate lists, candidate manifestos, the ballot casting interface, and the administrator election management console.
    \item \textbf{Application Relayer Tier (Node.js \& Express):} A secure backend server that authenticates incoming HTTP requests via JWT, enforces role-based access control, executes atomic database claims, and signs Hedera consensus transactions using the operator private key.
    \item \textbf{State Caching Tier (MongoDB Community Server 8):} A high-speed document database used to maintain voter registration state, hashed credentials, candidate election manifestos, and local transaction records. It acts as the concurrency barrier that prevents double-voting before network dispatch.
    \item \textbf{Distributed Consensus Tier (Hedera HCS):} The decentralized trust engine that orders, timestamps, and immutably records all ballot submissions across the global Hashgraph network.
\end{itemize}

\subsection{Module Interoperability and Data Flow}
The operational lifecycle of an election ballot proceeds through eight sequential stages, depicted in Figure 3.2:

\begin{figure}[H]
    \centering
    \includegraphics[width=0.75\textwidth]{./media/image7.png}
    \caption{Complete End-to-End Operational Lifecycle Flowchart from Voter Authentication to Consensus Receipt}
    \label{fig:lifecycle}
\end{figure}

1. \textbf{Authentication:} The student logs in via the React frontend. The backend verifies credentials against bcrypt hashes and issues a JWT. \\
2. \textbf{Candidate Selection:} The voter browses registered candidates, reviews policy manifestos, and initiates ballot submission. \\
3. \textbf{Token Verification:} The Express server verifies the Bearer JWT and extracts the authenticated student ID. \\
4. \textbf{Atomic State Claim:} MongoDB executes \texttt{findOneAndUpdate(\{ voterId, hasVoted: false \})}. If already true, request aborts immediately. \\
5. \textbf{Payload Construction:} The relayer formats an anonymized payload containing candidate ID, timestamp, and metadata. \\
6. \textbf{Hashgraph Broadcast:} The relayer signs and submits a \texttt{TopicMessageSubmitTransaction} to Hedera Testnet. \\
7. \textbf{aBFT Finality:} Hedera consensus nodes order the message and generate a consensus receipt with nanosecond timestamp. \\
8. \textbf{Receipt Return:} The backend returns the transaction ID to the voter, who can independently look up the receipt on HashScan.

\begin{figure}[H]
    \centering
    \includegraphics[width=0.9\textwidth]{./media/image6.png}
    \caption{Cryptographic Relayer Sequence Diagram showing Client, Server, Database, and Ledger interactions}
    \label{fig:sequence}
\end{figure}

\begin{figure}[H]
    \centering
    \includegraphics[width=0.9\textwidth]{./media/image4.png}
    \caption{Decoupled Frontend and Backend Component Hierarchy}
    \label{fig:hierarchy}
\end{figure}

\subsection{Implementation Environment}

\subsubsection{Required System Tools}
The system requires the following runtime dependencies installed on host machines or execution containers:
\begin{itemize}[leftmargin=*]
    \item \textbf{Node.js Runtime:} Version 22.x LTS (JavaScript V8 engine supporting native fetch and ECMAScript 2024).
    \item \textbf{Container Runtime:} Docker Engine 26.x and Docker Compose v2 for multi-service container orchestration.
    \item \textbf{Database Engine:} MongoDB Community Server 8.0 with WiredTiger storage engine supporting document-level locking.
    \item \textbf{Hedera SDK:} \texttt{@hashgraph/sdk} v2.46.0 providing gRPC node transport and transaction serialization.
\end{itemize}

\subsubsection{Environment Variables}
All sensitive credentials, ports, and network targets are configured through environment variables loaded at startup. Table 3.2 details the required configuration parameters:

\begin{table}[H]
    \centering
    \caption{System Configuration and Environment Variables}
    \label{tab:env_vars}
    \begin{tabularx}{\textwidth}{@{}l X X@{}}
        \toprule
        \textbf{Variable Name} & \textbf{Typical Value} & \textbf{Description \& Security Role} \\
        \midrule
        \texttt{PORT} & 5000 & TCP port on which the Express REST API server listens. \\
        \texttt{MONGO\_URI} & mongodb://127.0.0.1:27017... & Connection URI for MongoDB database cluster. \\
        \texttt{JWT\_SECRET} & super\_secret\_jwt... & 256-bit symmetric secret used to sign and verify session tokens. \\
        \texttt{HEDERA\_NETWORK} & testnet & Target consensus network (testnet, mainnet, or localnet). \\
        \texttt{HEDERA\_ACCOUNT\_ID} & 0.0.9757391 & Hedera operator account ID funding consensus transaction fees. \\
        \texttt{HEDERA\_PRIVATE\_KEY} & 302e020100300506032b... & ED25519 or ECDSA operator private key used to sign transactions. \\
        \texttt{HCS\_TOPIC\_ID} & 0.0.10697996 & Public Hedera Consensus Service Topic ID dedicated to the election. \\
        \texttt{VITE\_BACKEND\_URL} & http://localhost:5000 & Public endpoint URL used by the React SPA to reach API routes. \\
        \bottomrule
    \end{tabularx}
\end{table}

\subsubsection{Key Packages \& Dependencies}
Table 3.3 summarizes the core software packages deployed across the backend relayer and frontend single-page application:

\begin{table}[H]
    \centering
    \caption{Key Packages \& Dependencies}
    \label{tab:packages}
    \begin{tabularx}{\textwidth}{@{}l c l X@{}}
        \toprule
        \textbf{Package Name} & \textbf{Version} & \textbf{Scope} & \textbf{Role \& Engineering Function} \\
        \midrule
        \texttt{express} & 4.19.2 & Backend & High-performance HTTP web server and REST routing engine. \\
        \texttt{@hashgraph/sdk} & 2.46.0 & Backend & Official Hedera SDK for cryptographic signing and HCS dispatch. \\
        \texttt{mongoose} & 8.3.1 & Backend & Object-Document Mapper (ODM) with schema validation and atomic operators. \\
        \texttt{jsonwebtoken} & 9.0.2 & Backend & Stateless session authentication and JWT payload verification. \\
        \texttt{bcrypt} & 5.1.1 & Backend & Cryptographic password hashing utilizing blowfish cipher with 10 salt rounds. \\
        \texttt{react} & 18.2.0 & Frontend & Component-based declarative user interface rendering library. \\
        \texttt{vite} & 5.2.0 & Frontend & Next-generation frontend build tooling and local development server. \\
        \texttt{axios} & 1.6.8 & Frontend & Promise-based HTTP client for dispatching asynchronous API requests. \\
        \bottomrule
    \end{tabularx}
\end{table}

\subsubsection{Runtime Considerations}
Operating an electronic voting system under production election traffic requires strict adherence to runtime constraints:
\begin{itemize}[leftmargin=*]
    \item \textbf{Non-Blocking Asynchronous I/O:} The Node.js event loop handles thousands of concurrent HTTP connections by offloading crypto operations (bcrypt) and Hedera network gRPC calls to the internal libuv worker thread pool, preventing thread starvation.
    \item \textbf{Hedera Precheck \& Node Retries:} The Hedera SDK executes two-phase validation: local node prechecks (verifying signature format, operator account balance, and transaction syntax) followed by network consensus. If an individual node is unresponsive, the SDK automatically cycles to another consensus node.
    \item \textbf{Database Indexing:} Unique sparse indexes on \texttt{voterId} (\texttt{\{ unique: true \}}) in MongoDB ensure $O(1)$ query complexity during the atomic claim step, eliminating table scans during high-traffic voting spikes.
\end{itemize}

\subsubsection{Logging \& Artifact Locations}
To enable comprehensive post-election audits and operational observability, the platform maintains structured logging across all system tiers, illustrated in Figure 3.5:

\begin{figure}[H]
    \centering
    \includegraphics[width=0.85\textwidth]{./media/image3.png}
    \caption{Centralized Logging, Observability \& Mirror Node Pipeline}
    \label{fig:logging}
\end{figure}

\begin{itemize}[leftmargin=*]
    \item \textbf{Console \& Application Logs:} The backend server logs all incoming API requests, authenticated voter IDs, MongoDB query latencies, and Hedera transaction IDs.
    \item \textbf{Consensus Receipts (DLT Artifacts):} Every successful vote produces a Hedera Transaction Receipt containing the transaction ID (e.g., \texttt{0.0.9757391@1730182450.123456789}) and the consensus timestamp, returned directly to the client as an auditable receipt.
    \item \textbf{Public Mirror Node Archives:} Hedera mirror nodes archive all raw messages submitted to Topic \texttt{0.0.10697996}. These records are permanently queryable via REST.
\end{itemize}

% ------------------------------------------
% SECTION 4
% ------------------------------------------
\section{Implementation and Details}

\subsection{Orchestration and Entry Points}
The platform contains three primary execution entry points:
\begin{itemize}[leftmargin=*]
    \item \textbf{Topic Initialization Script (\texttt{createTopic.js}):} A standalone Node.js administrative script executed prior to election commencement. It connects to Hedera Testnet via the operator key, executes a \texttt{TopicCreateTransaction}, assigns the topic memo ('Hybrid Cloud E-Voting System'), and outputs the resulting HCS Topic ID for inclusion in \texttt{backend/.env}.
    \item \textbf{Backend API Server (\texttt{server.js}):} The primary backend process. It establishes database connections, seeds demonstration candidate data, initializes the Hedera Client, sets up CORS and JSON body parsing, and mounts authentication and election API routes.
    \item \textbf{Voter Web Interface (\texttt{Login.jsx \& App.jsx}):} The user-facing client application. It manages local JWT storage, executes role-based view routing, presents interactive ballots, and handles real-time toast notifications for system feedback.
\end{itemize}

\subsection{Authentication and Role-Based Access Control}
User authentication and route protection are handled via a decoupled JWT service. Figure 4.1 shows the credential verification and token issuance flow:

\begin{figure}[H]
    \centering
    \includegraphics[width=0.6\textwidth]{./media/image5.png}
    \caption{Authentication Flow \& JWT Claim Issuance for Voters and Administrators}
    \label{fig:auth_flow}
\end{figure}

When a user submits credentials, the server sanitizes the voter identifier, performs a lookup in the Voters collection, and compares the plaintext password against the stored bcrypt hash using \texttt{bcrypt.compare()}. Upon verification, the server generates a signed JWT containing the voter's ID (\texttt{sub}) and assigned role (Admins or Voters) with a 24-hour expiration window.

Subsequent requests to protected routes pass through two middleware functions:
\begin{itemize}
    \item \texttt{authenticateToken}: Extracts the Bearer token from the HTTP Authorization header, verifies signature validity via \texttt{jwt.verify()}, and attaches \texttt{req.voterId} and \texttt{req.role} to the request object.
    \item \texttt{requireRole(role)}: Inspects \texttt{req.role} and rejects non-authorized requests with HTTP 403 Forbidden, preventing regular voters from accessing administrative election management endpoints.
\end{itemize}

\subsection{Operator-Relayer Voting Engine \& Concurrency Control}
To allow voters to participate without personal Web3 wallets or gas fees, the backend implements an operator-relayer pattern. The server maintains a funded Hedera account. When an authenticated voter submits a ballot selection, the relayer constructs the transaction, signs it with the operator private key, and pays the fixed \$0.0001 network fee on behalf of the voter.

\subsection{The Atomic Ballot Claim \& Ledger Rollback Algorithm}
The primary challenge in electronic voting under high concurrency is preventing double-voting. If two requests for the same voter are processed simultaneously, a naive check-then-update pattern ('read \texttt{hasVoted} $\rightarrow$ submit to ledger $\rightarrow$ set \texttt{hasVoted = true}') creates a race window where both requests see \texttt{hasVoted == false} and both submit ballots to the ledger.

Our system eliminates this vulnerability by executing an atomic state claim inside MongoDB before dispatching the message to Hedera. Furthermore, if the consensus network call fails (e.g., due to network disconnection or timeout), an atomic rollback resets the voter's record, ensuring they are not unfairly disenfranchised:

\begin{lstlisting}[language=JavaScript, caption={Atomic Ballot Claim \& Ledger Rollback Implementation (backend/server.js)}]
// Step 1: Atomic database claim - prevents race conditions
const voter = await Voter.findOneAndUpdate(
  { voterId: req.voterId, hasVoted: false },
  { hasVoted: true, votedAt: new Date() },
  { returnDocument: 'after' }
);

if (!voter) {
  return res.status(403).json({ error: 'This voter has already voted.' });
}

// Step 2: Submit vote payload to Hedera Consensus Service
try {
  const payload = JSON.stringify({
      candidateId,
      timestamp: new Date().toISOString()
  });

  const tx = await new TopicMessageSubmitTransaction({
      topicId,
      message: payload
  }).execute(client);

  const receipt = await tx.getReceipt(client);

  res.json({
      message: 'Vote recorded on Hedera!',
      txId: tx.transactionId.toString(),
      consensusTimestamp: receipt.consensusTimestamp
  });

} catch (hederaErr) {
  // Step 3: Rollback voter claim if ledger broadcast fails
  await Voter.findOneAndUpdate(
      { voterId: req.voterId },
      { hasVoted: false, votedAt: null }
  );

  res.status(502).json({ error: 'Ledger broadcast failed. Voter status rolled back.' });
}
\end{lstlisting}

\subsection{Mirror Node Public Audit \& Explorer Verification}
To achieve zero-trust public verifiability, the election tally is not calculated from our internal MongoDB database. Instead, any auditor can query public Hedera Mirror Nodes directly via HTTP GET requests:
\texttt{https://testnet.mirrornode.hedera.com/api/v1/topics/0.0.10697996/messages}

The mirror node returns an array of messages ordered by consensus timestamp. The auditor base64-decodes the message payload, extracts the \texttt{candidateId}, and tallies the votes independently. Because the ledger is immutable and maintained by global consensus nodes, no centralized administrator can alter the final result.

% ------------------------------------------
% SECTION 5
% ------------------------------------------
\section{Results and Analysis}

\subsection{Qualitative Results (Artifacts \& Verification)}
The functional correctness of the system was validated through an end-to-end election simulation. Figures D.1 through D.8 (detailed in Appendix D) capture key execution milestones:
\begin{itemize}[leftmargin=*]
    \item \textbf{Role-Based Login:} Secure authentication interface offering instant demo credential autofill for testing.
    \item \textbf{Voter Workspace:} Live display of active candidates, voting history, and campus issue reporting.
    \item \textbf{Manifesto Tracking:} Voters review candidate campaign promises and track historical performance.
    \item \textbf{Cryptographic Confirmation:} Modal dialog prompting final ballot confirmation prior to Hedera submission.
    \item \textbf{Consensus Receipt:} Voters receive a unique transaction ID verifiable on the public HashScan block explorer.
\end{itemize}

\subsection{Quantitative Performance Metrics}
Performance benchmarks were conducted by executing concurrent voter load tests against the backend relayer connected to the live Hedera Testnet. Table 5.1 summarizes system performance across varying concurrency levels:

\begin{table}[H]
    \centering
    \caption{Concurrency Benchmark and Reliability Metrics}
    \label{tab:concurrency}
    \begin{tabularx}{\textwidth}{@{}l c c X X@{}}
        \toprule
        \textbf{Concurrent Voters} & \textbf{Success Rate} & \textbf{Avg Finality} & \textbf{Duplicate Rejections} & \textbf{Ledger Consistency} \\
        \midrule
        5 Voters & 100\% & 3.8 seconds & 0 attempted & 100\% Consistent \\
        10 Voters & 100\% & 4.1 seconds & 0 attempted & 100\% Consistent \\
        25 (Simulated Race) & 100\% & 4.5 seconds & 100\% Rejected (403) & Zero Double-Votes \\
        50 (Burst Load) & 98\% & 4.8 seconds & 100\% Rejected (403) & Ledger Matches DB \\
        \bottomrule
    \end{tabularx}
\end{table}

The system maintained an average consensus latency between 3.8 and 4.8 seconds across all test batches. Under race condition testing (where identical voter tokens submitted parallel requests simultaneously), MongoDB's atomic \texttt{findOneAndUpdate} operator successfully rejected 100\% of duplicate attempts with HTTP 403 Forbidden, proving total immunity to double-voting exploits.

\begin{table}[H]
    \centering
    \caption{Formal Threat Model Evaluation (STRIDE Framework)}
    \label{tab:stride}
    \begin{tabularx}{\textwidth}{@{}l X X l@{}}
        \toprule
        \textbf{Threat Category} & \textbf{Attack Vector Evaluated} & \textbf{Architectural Mitigation} & \textbf{Status} \\
        \midrule
        Spoofing & Credential stuffing \& impersonation & Bcrypt salt rounds (10), JWT expiration (24h) & Verified \\
        Tampering & Direct database tally manipulation & Decoupled HCS ledger anchoring with aBFT & Verified \\
        Repudiation & Voter claiming ballot was not recorded & Cryptographic transaction receipt \& HashScan URL & Verified \\
        Information Disclosure & Exposing voter candidate choices & Ballots decoupled from voter PII in ledger payloads & Verified \\
        Denial of Service & High-frequency concurrent ballot flooding & Atomic MongoDB \texttt{findOneAndUpdate} \& token rate limiting & Verified \\
        Elevation of Privilege & Voter accessing election admin routes & \texttt{requireRole('Admins')} middleware enforcement & Verified \\
        \bottomrule
    \end{tabularx}
\end{table}

\begin{table}[H]
    \centering
    \caption{OWASP Top 10 Security Assessment}
    \label{tab:owasp}
    \begin{tabularx}{\textwidth}{@{}l X X@{}}
        \toprule
        \textbf{OWASP Category} & \textbf{Vulnerability Risk} & \textbf{Engineering Control Implemented} \\
        \midrule
        A01: Broken Access Control & Unauthorized ballot submission & JWT Bearer token verification and \texttt{requireRole} middleware. \\
        A02: Cryptographic Failures & Voter password theft via database dump & Bcrypt salted one-way hashing; TLS/HTTPS in transit. \\
        A03: Injection & NoSQL injection into login queries & Strict string sanitization (\texttt{.trim().toUpperCase()}) and Mongoose schemas. \\
        A04: Insecure Design & Double-voting under distributed race & Atomic \texttt{findOneAndUpdate} state claim before ledger broadcast. \\
        A05: Security Misconfig & Exposing secrets in repository & Centralized \texttt{.env} configuration excluded from version control. \\
        A07: Identification Failures & Session hijacking & Stateless 24-hour expiration tokens signed with 256-bit secret. \\
        \bottomrule
    \end{tabularx}
\end{table}

\begin{table}[H]
    \centering
    \caption{Automated Test Suite Execution Results}
    \label{tab:tests}
    \begin{tabularx}{\textwidth}{@{}l X X l@{}}
        \toprule
        \textbf{Test Suite File} & \textbf{Component Under Test} & \textbf{Specific Scenarios Validated} & \textbf{Result} \\
        \midrule
        \texttt{tests/auth.test.js} & Authentication Controller & Credential rejection, password checks, token signing & PASS (5/5) \\
        \texttt{tests/hederaErrors.js} & Hedera Error Translator & \texttt{INVALID\_SIGNATURE} translation, network precheck retries & PASS (2/2) \\
        \textbf{Total Suite} & \textbf{Core Relayer \& DLT Bridge} & 7 automated regression tests verifying security boundaries & \textbf{PASS (100\%)} \\
        \bottomrule
    \end{tabularx}
\end{table}

\subsection{Edge Cases and Limitations}
Four critical edge cases were identified and addressed during evaluation:
\begin{itemize}[leftmargin=*]
    \item \textbf{Network Disconnection During Consensus Submission:} If a socket error occurs after MongoDB claims \texttt{hasVoted: true} but before Hedera returns a receipt, the catch block executes an atomic rollback setting \texttt{hasVoted: false}, allowing the student to re-cast their vote.
    \item \textbf{Operator Account Balance Depletion:} Because the relayer sponsors all transactions, if the operator account exhausts its HBAR balance, subsequent votes receive an \texttt{INSUFFICIENT\_PAYER\_BALANCE} error. In production, automated balance monitoring triggers automated re-funding alerts.
    \item \textbf{Mirror Node Eventual Consistency:} Public mirror nodes index consensus messages asynchronously, introducing a 1 to 2-second delay between consensus finality and availability on the REST API. The frontend displays the local transaction receipt immediately while notifying the user that public explorer indexing is processing.
    \item \textbf{Replay Attacks:} Submitting intercepted vote payloads to the backend is rejected because each JWT can only be used once per election cycle, as enforced by the database claim.
\end{itemize}

\subsection{Log Analysis \& Observability}
Examination of server and mirror node logs during load testing validated the strict ordering of consensus events. Each transaction receipt contains a sequence number that increments monotonically (e.g., Sequence \#1, \#2, \#3), guaranteeing that consensus timestamp order is globally unique and tamper-evident across the distributed network.

% ------------------------------------------
% SECTION 6
% ------------------------------------------
\section{Conclusion and Future Work}

\subsection{Conclusion}
This project successfully designed, implemented, and experimentally evaluated a hybrid Web2.5 electronic voting platform combining the accessibility of modern web technologies with the tamper resistance of the Hedera Consensus Service. By replacing centralized database authority with decentralized Hashgraph consensus, the platform guarantees that election results cannot be modified by any single administrative party. Concurrently, the operator-relayer pattern eliminates the complexities of cryptocurrency gas fees and Web3 browser wallets, enabling ordinary voters to participate seamlessly.

Quantitative benchmarks demonstrate that the platform achieves consensus finality within 3.8 to 4.8 seconds, maintains 100\% prevention of duplicate voting under high-concurrency race conditions, and provides end-to-end auditability via public mirror nodes and HashScan explorer receipts at a fixed transaction cost of \$0.0001 per vote. The platform validates that hybrid distributed ledger architectures provide a practical, highly secure foundation for institutional, collegiate, and municipal elections.

\subsection{Future Enhancements}
Future iterations will extend the platform in three strategic directions:
\begin{itemize}[leftmargin=*]
    \item \textbf{Zero-Knowledge Proofs (zk-SNARKs):} Integrating client-side zero-knowledge proofs will enable voters to prove their electoral eligibility and ballot validity without revealing their candidate selection or voter ID to the relayer backend, ensuring complete cryptographic ballot secrecy.
    \item \textbf{W3C Decentralized Identifiers (DIDs):} Replacing username/password authentication with Verifiable Credentials issued by university identity systems will establish self-sovereign digital identity for all academic participants.
    \item \textbf{Multi-Signature Operator Governance:} Transitioning the backend operator key to a threshold multi-signature account requiring consensus among multiple independent faculty trustees before election topics can be initialized or closed.
\end{itemize}

% ------------------------------------------
% REFERENCES
% ------------------------------------------
\newpage
\section{Supporting Materials / References}
\begin{enumerate}[label={[\arabic*]}, leftmargin=*]
    \item L. Baird, `The Swirlds Hashgraph Consensus Algorithm: Fair, Fast, Byzantine Fault Tolerance,' Swirlds Technical Report SWIRLDS-TR-2016-01, May 2016.
    \item Hedera Hashgraph, `Hedera Consensus Service (HCS) Whitepaper: Decentralized Trust for Enterprise Applications,' Hedera Documentation, 2020.
    \item B. Adida, `Helios: Web-based Open-Audit Voting,' in \textit{Proc. 17th USENIX Security Symposium}, San Jose, CA, 2008, pp. 335--348.
    \item P. McCorry, S. F. Shahandashti, and F. Hao, `A Smart Contract for Boardroom Voting with Decentralized Privacy,' in \textit{Financial Cryptography and Data Security}, LNCS vol. 10322, 2017, pp. 357--375.
    \item E. Androulaki et al., `Hyperledger Fabric: A Distributed Operating System for Permissioned Blockchains,' in \textit{Proc. 13th EuroSys Conference}, 2018, pp. 1--15.
    \item T. Kohno, A. Stubblefield, A. D. Rubin, and D. S. Wallach, `Analysis of an Electronic Voting System,' in \textit{IEEE Symposium on Security and Privacy}, 2004, pp. 27--40.
    \item A. J. Appel, `Security Flaws in an EVM Voting Machine,' in \textit{Proc. Electronic Voting Technology Workshop}, USENIX, 2011.
    \item D. Chaum, `Untraceable Electronic Mail, Return Addresses, and Digital Pseudonyms,' \textit{Communications of the ACM}, vol. 24, no. 2, pp. 84--90, 1981.
    \item D. Chaum et al., `Scantegrity II: End-to-End Verifiability for Optical Scan Election Systems,' \textit{IEEE Transactions on Information Forensics and Security}, vol. 4, no. 4, pp. 611--627, 2009.
    \item P. Y. Ryan et al., `Prêt à Voter: A Voter-Verifiable Voting System,' \textit{IEEE Transactions on Information Forensics and Security}, vol. 4, no. 4, pp. 662--673, 2009.
    \item R. Haenni et al., `A General Framework for Verifiable Online Voting,' Swiss Federal Chancellery Technical Report, 2017.
    \item S. Springall et al., `Security Analysis of the Estonian Internet Voting System,' in \textit{ACM SIGSAC Conference on Computer and Communications Security}, 2014, pp. 703--715.
    \item S. Heiberg and J. Willemson, `Verifiable Internet Voting in Estonia,' in \textit{IEEE 3rd International Conference on e-Voting and Identity}, 2014.
    \item M. B. Jones, J. Bradley, and N. Sakimura, `JSON Web Token (JWT),' RFC 7519, Internet Engineering Task Force (IETF), May 2015.
    \item Open Web Application Security Project (OWASP), `OWASP Top Ten Web Application Security Risks,' OWASP Foundation, 2021.
    \item National Institute of Standards and Technology (NIST), `Digital Identity Guidelines: Authentication and Lifecycle Management,' NIST SP 800-63B, 2020.
    \item React Documentation, `React v18.2 Component Lifecycle and Hooks Architecture,' Meta Open Source, 2024.
    \item MongoDB Inc., `MongoDB 8.0 Manual: Atomic Transactions and Concurrency Control,' MongoDB Documentation, 2024.
    \item Amazon Web Services, `AWS Serverless Architecture: API Gateway, AWS Lambda, and Amazon Cognito Reference Architecture,' AWS Whitepapers, 2023.
    \item Hedera Documentation, `Hedera Mirror Node Architecture \& REST API Guide,' Hedera Developer Portal, 2026.
    \item HashScan, `Hedera Network Ledger Explorer,' Available: \url{https://hashscan.io}, 2026.
    \item Node.js Foundation, `Node.js v22 LTS Core API Specification,' OpenJS Foundation, 2024.
\end{enumerate}

% ------------------------------------------
% APPENDICES
% ------------------------------------------
\newpage
\appendix
\section{Appendix}

\subsection{Appendix A: Full Code Listing of Key Modules}
This appendix presents the complete, verified source code listings for the core modules supporting the voting platform.

\begin{lstlisting}[language=JavaScript, caption={Backend Relayer \& Consensus Engine (backend/server.js)}]
const express = require("express");
const cors = require("cors");
const mongoose = require("mongoose");
const jwt = require("jsonwebtoken");
const bcrypt = require("bcrypt");
const { Client, PrivateKey, TopicMessageSubmitTransaction } = require("@hashgraph/sdk");
const path = require("path");
require("dotenv").config({ path: path.resolve(__dirname, ".env") });

const Voter = require("./models/Voter");
const Candidate = require("./models/Candidate");
const Election = require("./models/Election");
const CampusIssue = require("./models/CampusIssue");
const EvidenceDoc = require("./models/EvidenceDoc");
const { seedManifestoData } = require("./seedData");
const { login } = require("./controllers/auth");
const { getVoteBroadcastErrorMessage } = require("./lib/voteErrors");

// Environment validation
const requiredEnvVars = [
  'MONGO_URI',
  'HEDERA_NETWORK',
  'HEDERA_ACCOUNT_ID',
  'HEDERA_PRIVATE_KEY',
  'HCS_TOPIC_ID',
  'JWT_SECRET'
];

for (const envVar of requiredEnvVars) {
  if (!process.env[envVar]) {
    console.error(`FATAL: Missing required environment variable: ${envVar}`);
    process.exit(1);
  }
}

const app = express();
app.use(cors());
app.use(express.json());

const JWT_SECRET = process.env.JWT_SECRET;

// MongoDB Connection
mongoose.connect(process.env.MONGO_URI)
  .then(async () => {
    console.log("Connected to MongoDB.");
    // Seed manifesto demo data if no candidates exist yet
    try {
      const candidateCount = await Candidate.countDocuments();
      if (candidateCount === 0) {
        console.log("No candidates found. Seeding manifesto demonstration data...");
        await seedManifestoData();
        console.log("Manifesto seed data loaded successfully.");
      } else {
        console.log(`Database already contains ${candidateCount} candidate(s). Skipping seed.`);
      }
    } catch (seedErr) {
      console.error("Manifesto seed warning (non-fatal):", seedErr.message);
    }
  })
  .catch(err => console.error("MongoDB Connection Error:", err));

// Hedera Client Initialization
let client;
if (process.env.HEDERA_NETWORK === "localnet") {
  client = Client.forLocalnet();
} else {
  client = Client.forTestnet();
}
const hederaAccountId = process.env.HEDERA_ACCOUNT_ID.trim();
const hederaPrivateKey = process.env.HEDERA_PRIVATE_KEY.trim();

client.setOperator(
  hederaAccountId,
  hederaPrivateKey.startsWith('0x') || hederaPrivateKey.length === 66
    ? PrivateKey.fromStringECDSA(hederaPrivateKey)
    : PrivateKey.fromString(hederaPrivateKey)
);
const topicId = process.env.HCS_TOPIC_ID;

// Authentication middleware
const authenticateToken = (req, res, next) => {
  const authHeader = req.headers['authorization'];
  const token = authHeader && authHeader.split(' ')[1];

  if (!token) return res.status(401).json({ error: 'Access token required' });

  jwt.verify(token, JWT_SECRET, (err, user) => {
    if (err) return res.status(403).json({ error: 'Invalid or expired token' });

    // Set voterId and role from token claims
    req.voterId = user.sub;
    req.role = user.role;
    next();
  });
};

// Role-based middleware
const requireRole = (requiredRole) => {
  return (req, res, next) => {
    if (!req.voterId || !req.role) {
      return res.status(403).json({ error: 'Insufficient permissions' });
    }

    if (req.role !== requiredRole) {
      return res.status(403).json({ error: 'Insufficient permissions' });
    }
    next();
  };
};

app.post('/api/auth/login', login);

// Voter: Submit Ballot to Hedera HCS
app.post("/castVote", authenticateToken, requireRole("Voters"), async (req, res) => {
  try {
    const election = await Election.findOne();
    if (!election || !election.isActive) {
      return res.status(403).json({ error: "Voting is currently closed." });
    }

    const voterId = req.voterId; // Extract from JWT sub claim
    const { candidateId } = req.body;

    // Validate candidate
    const candidate = await Candidate.findOne({ candidateId });
    if (!candidate) {
      return res.status(400).json({ error: "Invalid candidate selected." });
    }

    // Atomically claim the voter's ballot before broadcasting it.
    const voter = await Voter.findOneAndUpdate(
      { voterId: voterId, hasVoted: false },
      { hasVoted: true, votedAt: new Date() },
      { returnDocument: 'after' }
    );

    if (!voter) {
      const existingVoter = await Voter.findOne({ voterId: voterId });
      if (!existingVoter) {
        return res.status(401).json({ error: "Voter ID not found in electoral roll." });
      } else {
        return res.status(403).json({ error: "This voter has already voted." });
      }
    }

    // Submit payload to Hedera Consensus Service
    const payload = JSON.stringify({
      candidateId,
      timestamp: new Date().toISOString()
    });

    let txResponse;
    let receipt;
    try {
      txResponse = await new TopicMessageSubmitTransaction({
        topicId,
        message: payload,
      }).execute(client);
      receipt = await txResponse.getReceipt(client);
    } catch (hederaErr) {
      // Revert voter update on Hedera failure
      await Voter.findOneAndUpdate(
        { voterId: voterId },
        { hasVoted: false, votedAt: null }
      );
      console.error("Hedera transaction failed, voter rollback executed:", hederaErr);
      return res.status(502).json({ error: getVoteBroadcastErrorMessage(hederaErr) });
    }

    res.json({
      message: "Vote recorded on Hedera!",
      transactionId: txResponse.transactionId.toString(),
      status: receipt.status.toString()
    });
  } catch (err) {
    console.error("Ballot submission error:", err);
    res.status(500).json({ error: "Failed to broadcast vote to ledger." });
  }
});

const PORT = process.env.PORT || 5000;
app.listen(PORT, () => console.log(`Server running on port ${PORT}`));
\end{lstlisting}

\newpage
\begin{lstlisting}[language=JavaScript, caption={Hedera Topic Creation Script (backend/createTopic.js)}]
require("dotenv").config();
const { Client, PrivateKey, TopicCreateTransaction } = require("@hashgraph/sdk");

async function main() {
    const myAccountId = process.env.HEDERA_ACCOUNT_ID;
    const myPrivateKey = PrivateKey.fromString(process.env.HEDERA_PRIVATE_KEY);

    const client = Client.forTestnet();
    client.setOperator(myAccountId, myPrivateKey);

    console.log("Creating new HCS Topic...");

    const transaction = new TopicCreateTransaction().setTopicMemo("Hybrid Cloud E-Voting System");
    const txResponse = await transaction.execute(client);

    const receipt = await txResponse.getReceipt(client);
    const topicId = receipt.topicId;

    console.log(`\n[SUCCESS] Your New Topic ID is: ${topicId.toString()}`);
    console.log("Copy this ID and add it to your .env file as HCS_TOPIC_ID\n");

    process.exit();
}

main().catch((err) => {
    console.error(err);
});
\end{lstlisting}

\newpage
\begin{lstlisting}[language=JavaScript, caption={Authentication Controller (backend/controllers/auth.js)}]
const jwt = require("jsonwebtoken");
const bcrypt = require("bcrypt");
const Voter = require("../models/Voter");

const JWT_SECRET = process.env.JWT_SECRET;

const login = async (req, res) => {
  try {
    const { voterId, password } = req.body;

    if (!voterId || !password) {
      return res.status(400).json({ error: "Voter ID and password are required" });
    }

    const cleanVoterId = voterId.trim().toUpperCase();
    const voter = await Voter.findOne({ voterId: cleanVoterId });

    if (!voter) {
      return res.status(401).json({ error: "Invalid credentials" });
    }

    const isMatch = await bcrypt.compare(password, voter.password);
    if (!isMatch) {
      return res.status(401).json({ error: "Invalid credentials" });
    }

    // Determine role for token
    const role = voter.role === "admin" ? "Admins" : "Voters";

    // Generate JWT token
    const token = jwt.sign(
      {
        sub: voter.voterId,
        role: role
      },
      JWT_SECRET,
      { expiresIn: "24h" }
    );

    res.json({
      message: "Login successful",
      token,
      user: {
        voterId: voter.voterId,
        role: role
      }
    });
  } catch (err) {
    console.error("Login error:", err);
    res.status(500).json({ error: "Internal server error" });
  }
};

module.exports = { login };
\end{lstlisting}

\newpage
\subsection{Appendix B: Deployment \& Docker Configuration}

\begin{lstlisting}[language=bash, caption={Environment Variable Specification (backend/.env)}]
# Server Configuration
PORT=5000

# Database Connection
MONGO_URI=mongodb://127.0.0.1:27017/voting_system_db

# Authentication Secrets
JWT_SECRET=super_secret_jwt_voting_key_2026
ADMIN_PASSWORD=AdminPassword123!

# Hedera Network & Credentials
HEDERA_NETWORK=testnet
HEDERA_ACCOUNT_ID=0.0.9757391
HEDERA_PRIVATE_KEY=302e020100300506032b6570042204201da2a618d47ba40c96e303afca89c3adcd9a2c8da303173796b5af0df8477889

# Hedera Consensus Service (HCS) Topic ID
HCS_TOPIC_ID=0.0.10697996
\end{lstlisting}

\begin{lstlisting}[language=bash, caption={Multi-Container Deployment Specification (docker-compose.yml)}]
services:
  mongodb:
    image: mongo:8
    container_name: voting-mongodb
    ports:
      - "27017:27017"
    volumes:
      - mongodb_data:/data/db
    environment:
      - MONGO_INITDB_DATABASE=evoting_db
    healthcheck:
      test: ["CMD", "mongosh", "--quiet", "--eval", "db.adminCommand('ping').ok"]
      interval: 5s
      timeout: 5s
      retries: 10

  backend:
    build: ./backend
    container_name: voting-backend
    ports:
      - "5000:5000"
    env_file:
      - ./backend/.env
    environment:
      MONGO_URI: mongodb://mongodb:27017/evoting_db
      PORT: 5000
      BACKEND_URL: http://localhost:5000
    depends_on:
      mongodb:
        condition: service_healthy

  frontend:
    build:
      context: ./frontend
      args:
        VITE_BACKEND_URL: ${VITE_BACKEND_URL:-http://localhost:5000}
    container_name: voting-frontend
    ports:
      - "5173:80"
    depends_on:
      - backend

volumes:
  mongodb_data:
\end{lstlisting}

\newpage
\subsection{Appendix C: User Manual and Operator Guide}

\subsubsection{C.1 System Requirements \& Installation}
\begin{enumerate}[leftmargin=*]
    \item Clone the project repository to your local workstation.
    \item Ensure Node.js (v20+ or v22 LTS) and MongoDB are installed and active.
    \item Navigate to the backend directory and install dependencies:
    \begin{verbatim}
    cd backend && npm install
    \end{verbatim}
    \item Configure environment credentials in \texttt{backend/.env} with your Hedera Testnet Account ID and Private Key.
    \item Initialize a new Hedera Consensus Topic:
    \begin{verbatim}
    node createTopic.js
    \end{verbatim}
    Copy the generated Topic ID (e.g., 0.0.10697996) and assign it to \texttt{HCS\_TOPIC\_ID} in \texttt{.env}.
    \item Start the backend server:
    \begin{verbatim}
    npm start
    \end{verbatim}
    \item Navigate to the frontend directory and start the Vite development server:
    \begin{verbatim}
    cd ../frontend && npm install && npm run dev
    \end{verbatim}
    \item Access the web portal at \texttt{http://localhost:5173} (or \texttt{http://localhost:80} via Docker).
\end{enumerate}

\subsubsection{C.2 Voter Guide: Casting a Verified Ballot}
\begin{enumerate}[leftmargin=*]
    \item Open the application in any modern web browser.
    \item On the login screen, select 'Login as Voter'. For demo testing, click the 'Voter Demo' button to autofill credentials.
    \item Upon authentication, browse the active candidates and examine their campaign manifestos.
    \item Click 'Vote' next to your preferred candidate. A confirmation modal will appear.
    \item Confirm your selection. The backend will atomically claim your ballot and broadcast it to Hedera Consensus Service.
    \item Upon completion, copy your Transaction ID and verify your ballot on the public HashScan explorer.
\end{enumerate}

\subsubsection{C.3 Auditor Guide: Public Ledger Verification on HashScan}
\begin{enumerate}[leftmargin=*]
    \item Open \texttt{https://hashscan.io/testnet} in a browser.
    \item Enter Topic ID 0.0.10697996 in the search bar.
    \item Review the ordered list of messages with nanosecond timestamps.
    \item Query the public mirror node endpoint via \texttt{curl} or Postman:
    \begin{verbatim}
    curl https://testnet.mirrornode.hedera.com/api/v1/topics/0.0.10697996/messages
    \end{verbatim}
    \item Base64-decode each message payload to independently sum all candidate selections and verify the official election tally.
\end{enumerate}

\newpage
\subsection{Appendix D: Sample Output Visuals (Figures D.1 -- D.8)}
This appendix documents the production visual artifacts produced during a standard election cycle.

\begin{figure}[H]
    \centering
    \includegraphics[width=0.9\textwidth]{./media/image8.jpeg}
    \caption{Dual-Role Authentication Interface with Voter/Admin demo fast-login buttons}
    \label{fig:d1}
\end{figure}

\begin{figure}[H]
    \centering
    \includegraphics[width=0.9\textwidth]{./media/image9.jpeg}
    \caption{Live Voter Dashboard showing election title, active candidates, and issue reporting tabs}
    \label{fig:d2}
\end{figure}

\begin{figure}[H]
    \centering
    \includegraphics[width=0.9\textwidth]{./media/image10.jpeg}
    \caption{Candidate Manifesto \& Track Record Verification Interface}
    \label{fig:d3}
\end{figure}

\begin{figure}[H]
    \centering
    \includegraphics[width=0.9\textwidth]{./media/image11.jpeg}
    \caption{Administrative Election Management Console with live election controls}
    \label{fig:d4}
\end{figure}

\begin{figure}[H]
    \centering
    \includegraphics[width=0.9\textwidth]{./media/image12.jpeg}
    \caption{Candidate Registration \& Manifesto Management Portal}
    \label{fig:d5}
\end{figure}

\begin{figure}[H]
    \centering
    \includegraphics[width=0.9\textwidth]{./media/image13.jpeg}
    \caption{Cryptographic Ballot Selection \& Confirmation Modal}
    \label{fig:d6}
\end{figure}

\begin{figure}[H]
    \centering
    \includegraphics[width=0.9\textwidth]{./media/image14.jpeg}
    \caption{Public HashScan Ledger Consensus Receipt showing confirmed Transaction ID}
    \label{fig:d7}
\end{figure}

\begin{figure}[H]
    \centering
    \includegraphics[width=0.9\textwidth]{./media/image15.jpeg}
    \caption{Live Hedera Mirror Node JSON Message Snapshot verifying consensus sequencing}
    \label{fig:d8}
\end{figure}

\newpage
\subsection{Appendix E: Technical Acronyms and Formal Domain Glossary}
Table E.1 defines key architectural terms and distributed ledger concepts referenced throughout this report:

\begin{table}[H]
    \centering
    \caption{Technical Acronyms and Formal Domain Glossary}
    \label{tab:glossary}
    \begin{tabularx}{\textwidth}{@{}l X@{}}
        \toprule
        \textbf{Acronym / Term} & \textbf{Formal Engineering Definition \& Significance} \\
        \midrule
        \textbf{aBFT} & Asynchronous Byzantine Fault Tolerance: The highest security standard for consensus algorithms, guaranteeing network convergence despite arbitrary malicious delays or node compromises. \\
        \textbf{HCS} & Hedera Consensus Service: A decentralized service that orders and timestamps transactions onto the public Hashgraph ledger without smart contract gas auctions. \\
        \textbf{DLT} & Distributed Ledger Technology: A decentralized transactional database cryptographically synchronized across independent consensus nodes. \\
        \textbf{JWT} & JSON Web Token (RFC 7519): A compact, URL-safe means of representing claims securely between two parties, used for stateless authentication. \\
        \textbf{RBAC} & Role-Based Access Control: An access authorization mechanism that restricts API route execution based on assigned roles (e.g., Voters vs Admins). \\
        \textbf{SPA} & Single-Page Application: A modern web application that dynamically updates the browser DOM without full-page reloads, built using React and Vite. \\
        \textbf{REST} & Representational State Transfer: An architectural style for networked hypermedia applications communicating over standard HTTP methods. \\
        \textbf{Bcrypt} & A password-hashing function based on the Blowfish cipher incorporating a cryptographic salt to protect against rainbow table attacks. \\
        \textbf{TPS} & Transactions Per Second: The quantitative throughput measurement evaluating distributed ledger scalability. \\
        \textbf{Mirror Node} & An independent node in the Hedera network that archives historical consensus transactions and exposes public REST/gRPC query APIs. \\
        \textbf{HashScan} & The official public block explorer for Hedera networks, enabling transparent verification of transactions, topics, and account balances. \\
        \bottomrule
    \end{tabularx}
\end{table}

\end{document}
